\section{Conclusion}
\label{sec:conclusion}

This study demonstrates the effectiveness of multi-domain ViT representation fusion (Face $\oplus$ Emotion $\oplus$ Age) for the classification of six intersectional classes from three macro-racial groups (Asian, Black, White) on the DemogPairs dataset. Under the Strict Subject-Disjoint Partitioning protocol with 0\% identity leakage (480 training subjects and 120 held-out testing subjects), the LR classifier achieved peak performance with an accuracy of 88.66\%, macro precision of 88.75\%, macro recall of 88.66\%, and macro F1-score of 88.68\%, followed by the SVM classifier with an accuracy of 88.61\%, macro precision of 88.71\%, macro recall of 88.61\%, and macro F1-score of 88.63\%. Paired inferential statistical significance testing of McNemar demonstrated the significance of tri-domain fusion over the single-domain Face representation ($p < 0.001$) and confirmed predictive parity between LR and SVM ($p = 1.0000$, with a net difference of one sample out of 2,160 images). Algorithmic fairness evaluation across subgroups demonstrated prediction stability with a recall range from 84.72\% to 94.44\% and an EOD gap of 9.72~pp across both top-performing models. From the aspect of computational efficiency, the downstream LR classifier proved to be the most optimal with an inference latency of 0.0046~ms per sample and a model size of 0.11~MB, realizing an ultra-lightweight classification head for practical deployment.

Despite offering significant empirical contributions, this study has two primary limitations that warrant attention. First, the demographic scope of this investigation is limited to three macro-racial groups within the DemogPairs dataset (Asian, Black, White); thus, generalizing classification performance to multiracial groups and other ethnic populations still requires additional empirical testing. Second, the multi-domain representation extraction process is still conducted separately (offline) utilizing three frozen ViT-Base backbone models without unified end-to-end fine-tuning, thereby requiring a dedicated feature processing stage prior to downstream classification. In response to these limitations, future research directions are oriented toward three strategic priorities. First, cross-dataset validation needs to be conducted on large-scale external in-the-wild corpora, such as FairFace or UTKFace, to evaluate model robustness against extreme pose variations, occlusions, and dynamic illumination. Second, future research needs to develop an adaptive inter-domain cross-attention fusion mechanism (end-to-end cross-attention fusion) as well as apply quantization or knowledge distillation techniques. Third, integrating Explainable Artificial Intelligence (XAI) visualization methods, such as Attention Rollout, represents a crucial step to provide feature representation transparency and support Responsible AI principles.

\section*{Data Availability Statement}

The DemogPairs dataset analyzed in this study is publicly available for academic research purposes. The facial images along with demographic label annotations can be freely accessed via the developers' official publication at \url{https://ihupont.github.io/publications/2019-05-16-demogpairs} and the official Google Drive storage repository at \url{https://drive.google.com/file/d/1f_ez-ll6wxDXScrG4ceStZRuLT_K8Uy8/view?usp=sharing}.
